%%%%%%%%%%%%%%%%%%%%%%%%%%%%%%%%%%%%%%%%%%%%%%%%%%%%%%%%%%%%%%%%%%%%%%%%%%%%%%%%%%%%%%%%%%%%%%%%%%%%%%
%
%   Filename    : abstract-filipino.tex 
%
%   Description : This file will contain your abstract in filipino language.
%                 
%%%%%%%%%%%%%%%%%%%%%%%%%%%%%%%%%%%%%%%%%%%%%%%%%%%%%%%%%%%%%%%%%%%%%%%%%%%%%%%%%%%%%%%%%%%%%%%%%%%%%%

\renewcommand{\abstractname}{Abstrak}

\begin{abstract}

% Do not put citations or quotes in the abstract.

% Lubos na binago ng mga short-form video (SFV) platform gaya ng TikTok, YouTube Shorts, at Instagram Reels ang paglikha ng digital na nilalaman. Gayunman, nahaharap ang mga content creator sa kawalang-katiyakan sa pagtaya ng magiging audience engagement bago mailathala ang kanilang mga video. Mahalaga ang maagang interaksiyon sapagkat ito ang kadalasang nagtatakda ng visibility ng nilalaman; subalit ang mga bagong upload na video ay kadalasang wala pang sapat na datos ng interaksiyon, kaya nagiging hamon ang paggawa ng desisyon bago ang publikasyon. Nakatuon ang mga naunang pag-aaral sa mga modelong platform-centered na gumagamit ng multimodal na katangian ng video—biswal, audio, at tekstuwal—ngunit hindi gaanong nabibigyang-pansin ang mga salik na may kaugnayan sa creator at sa konteksto. Layunin ng pag-aaral na ito na bumuo ng isang creator-oriented na modelo ng prediksyon ng engagement na pinagsasama ang mga multimodal na senyas ng nilalaman at mga katangian ng creator, tulad ng bilang ng tagasubaybay at kasaysayan ng pagpo-post, pati na ang mga kontekstuwal na salik gaya ng oras ng pagpo-post at mga padron ng aktibidad ng audience. Nilalayon ng modelo na matantiya ang potensiyal na engagement bago pa man i-post ang video upang masuportahan ang mga desisyong malikhain na nakabatay sa datos. Sa pag-uugnay ng nilalaman, creator, at konteksto, nag-aambag ang pananaliksik na ito sa human-centered na predictive analytics, sumusuporta sa estratehikong alokasyon ng yaman sa creator economy, at maaaring magtaguyod ng mas sari-sari at sustenableng ekosistem ng digital na nilalaman.

% Lubos na binago ng mga short-form video (SFV) platform gaya ng TikTok, YouTube Shorts, at Instagram Reels ang paglikha ng digital na nilalaman. Gayunman, nahaharap ang mga content creator sa kawalang-katiyakan sa pagtaya ng magiging audience engagement bago mailathala ang kanilang mga video. Mahalaga ang maagang interaksiyon sapagkat ito ang kadalasang nagtatakda ng visibility ng nilalaman; subalit ang mga bagong upload na video ay kadalasang wala pang sapat na datos ng interaksiyon, kaya nagiging hamon ang paggawa ng desisyon bago ang publikasyon. Nakatuon ang mga naunang pag-aaral sa mga modelong platform-centered na gumagamit ng multimodal na katangian ng video—biswal, audio, at tekstuwal—ngunit hindi gaanong nabibigyang-pansin ang mga salik na may kaugnayan sa creator at sa konteksto. Dahil dito, wala pang umiiral na modelo na direktang nagbibigay-daan sa mga indibidwal na creator na matantiya ang potensiyal na engagement ng kanilang sariling video bago ito i-upload. Layunin ng pag-aaral na ito na bumuo ng isang creator-oriented na modelo ng prediksyon ng engagement na pinagsasama ang mga multimodal na senyas ng nilalaman at mga katangian ng creator, tulad ng bilang ng tagasubaybay at kasaysayan ng pagpo-post, pati na ang mga kontekstuwal na salik gaya ng oras ng pagpo-post at mga padron ng aktibidad ng audience. Upang makamit ito, bubuo ang pag-aaral ng isang dataset sa pamamagitan ng boluntaryong data donation mula sa mga TikTok creator, kung saan pagsasamahin ang mga raw video file at opisyal na creator analytics. Gagamit ng ensemble ng mga Large Multimodal Model—ang VideoLLaMA2 at Qwen2.5-VL upang magsanib ang mga multimodal content embedding at ang structured na metadata ng creator at konteksto para sa pagtaya ng engagement . Nilalayon ng modelo na matantiya ang potensiyal na engagement bago pa man i-post ang video upang masuportahan ang mga desisyong malikhain na nakabatay sa datos. Sa pag-uugnay ng nilalaman, creator, at konteksto, nag-aambag ang pananaliksik na ito sa human-centered na predictive analytics, sumusuporta sa estratehikong alokasyon ng yaman sa creator economy, at maaaring magtaguyod ng mas sari-sari at sustenableng ekosistem ng digital na nilalaman.

Lubos na binago ng mga short-form video (SFV) platform gaya ng TikTok ang paglikha ng digital na nilalaman. Gayunman, nahaharap ang mga content creator sa kawalang-katiyakan sa pagtaya ng magiging audience engagement bago mailathala ang kanilang mga video. Mahalaga ang maagang interaksiyon sapagkat ito ang kadalasang nagtatakda ng visibility ng nilalaman; subalit ang mga bagong upload na video ay kadalasang wala pang sapat na datos ng interaksiyon, kaya nagiging hamon ang paggawa ng desisyon bago ang publikasyon. Nakatuon ang mga naunang pag-aaral sa mga modelong platform-centered na gumagamit ng multimodal na katangian ng video—biswal, audio, at tekstuwal—ngunit hindi gaanong nabibigyang-pansin ang mga salik na may kaugnayan sa creator at sa konteksto. Dahil dito, wala pang umiiral na bukas na dokumentadong modelo na komprehensibong nagbibigay-daan sa mga indibidwal na creator na matantiya ang potensiyal na engagement ng kanilang sariling video bago ito i-upload. Layunin ng pag-aaral na ito na bumuo ng isang creator-oriented na modelo ng prediksyon ng engagement na pinagsasama ang mga multimodal na senyas ng nilalaman at mga katangian ng creator, tulad ng bilang ng tagasubaybay at edad ng account, pati na ang mga kontekstuwal na salik gaya ng oras at araw ng pagpo-post. Nilalayon ng modelong ito na matantiya ang potensiyal na engagement—na isinasakatuparan bilang dalawang sukatan ng maagang viewer-retention, ang Engagement Continuation Rate (ECR) at Normalized Average Watch Percentage (NAWP)—bago pa man i-post ang video, upang masuportahan ang mga desisyong malikhain na nakabatay sa datos. Upang makamit ito, bubuo ang pag-aaral ng MicroTok-PH, isang dataset na bubuoin sa pamamagitan ng boluntaryong data donation mula sa mga Filipinong micro-influencer, kung saan pagsasamahin ang mga raw video file at opisyal na creator analytics. Ang iminumungkahing modelo, ang C3-LMM, ay gagamit ng isang frozen ensemble ng mga Large Multimodal Model—ang VideoLLaMA2 at Qwen2.5-VL—upang matantiya ang engagement bilang isang gawaing regression, sa pamamagitan ng dalawang estratehiya ng fusion na inihahambing ng pag-aaral: prompt-based feature injection at auxiliary-model ensemble. Sa pag-uugnay ng nilalaman, creator, at konteksto, mag-aambag ang pananaliksik na ito sa human-centered na predictive analytics, susuporta sa estratehikong alokasyon ng yaman sa creator economy, at maaaring magtaguyod ng mas sari-sari at sustenableng ekosistem ng digital na nilalaman.


\begin{flushleft}
\begin{tabular}{lp{4.25in}}
\hspace{-0.5em}\textbf{Keywords:}\hspace{0.25em} & short-form video, prediksyon ng engagement, pre-publication analytics, multimodal learning, creator-oriented na pagmomodelo
\end{tabular}
\end{flushleft}
\end{abstract}
